\section{Solving the Problem: Method}\label{sec:method}

\subsection{Why a first-order proximal method}

Problem \cref{eq:main} splits as $\min_w f(w) + g(w)$ with
\begin{equation}\label{eq:split}
f(w) = -\muhat^\top w + \kappa\norm{\Ssqrt w}_2 + \gamma\, w^\top \Sig w,
\qquad
g(w) = \lambda\norm{w}_1 + \mathbb{1}\{\one^\top w = 1\},
\end{equation}
where $\mathbb{1}\{\cdot\}$ is $0$ on the feasible set and $+\infty$ off it. The part $f$
admits a subgradient everywhere; $g$ is nonsmooth but, as \cref{prop:jointprox} shows, its
proximal operator can be computed \emph{exactly}. That is precisely the situation proximal
methods are built for \citep{parikh2014}. Folding the budget constraint into $g$ rather than
handling it separately is the decision that makes the method work, for reasons
\cref{sec:proxbug} makes concrete. The solver is pure NumPy: no \texttt{cvxpy},
\texttt{scipy.optimize}, or \texttt{sklearn} routine appears in it, only in the independent
verification of \cref{sec:insample}.

\subsection{The subgradient and step size}

At iterate $w_k$ the method takes a subgradient of $f$:
\begin{equation}\label{eq:subgrad}
v_k \;=\; -\muhat \;+\; 2\gamma\,\Sig w_k
        \;+\; \kappa\,\frac{\Sig w_k}{\norm{\Ssqrt w_k}_2}.
\end{equation}
The third term deserves derivation, being the step at which an implementation is most easily
and most silently wrong. With $u = \Ssqrt w$ and $\Ssqrt$ symmetric, the chain rule gives
\begin{equation}\label{eq:chainrule}
\nabla_w \norm{\Ssqrt w}_2
  \;=\; \bigl(\Ssqrt\bigr)^\top \frac{u}{\norm{u}_2}
  \;=\; \frac{\Ssqrt \Ssqrt w}{\norm{\Ssqrt w}_2}
  \;=\; \frac{\Sig w}{\norm{\Ssqrt w}_2}.
\end{equation}
The numerator is $\Sig w$, the \emph{full} covariance matrix --- not $\Ssqrt w$. Writing
$\Ssqrt w$ is a natural-looking substitution, since $\Ssqrt$ is what appears inside the norm,
and it yields a solver that runs without complaint and converges to the wrong point. Where
$\Ssqrt w = 0$, \cref{eq:chainrule} is undefined; the implementation takes the subgradient to
be $0$, legitimate because the subdifferential of $\norm{\cdot}_2$ at the origin is the closed
unit ball.

The descent step is $z_k = w_k - \alpha_k v_k$ with $\alpha_k = \alpha_0/\sqrt{k+1}$,
$\alpha_0 = 10$. This schedule is square-summable but not summable
($\sum_k\alpha_k = \infty$, $\sum_k\alpha_k^2 < \infty$), the standard sufficient condition for
subgradient-method convergence on a convex objective \citep{boyd2004}. A constant step would
drive the iterates only into a neighbourhood of the optimum whose radius scales with the step,
leaving an objective-gap floor no amount of iteration removes.

\subsection{The joint proximal step}

The remaining piece is the prox of $g$, stated as a proposition because getting it right rather
than approximately right is the technical core of the implementation.

\begin{proposition}[Exact prox of $\ell_1$ plus a budget constraint]\label{prop:jointprox}
Let $z \in \R^N$, let $t > 0$, and let
$\softt_t(x) = \operatorname{sign}(x)\max(|x| - t,\, 0)$ denote coordinate-wise soft
thresholding. The problem
\[
  \prox_{g}(z) \;=\; \arg\min_{w \in \R^N}\;
    \tfrac{1}{2}\norm{w - z}_2^2 + t\norm{w}_1
    \quad \text{subject to} \quad \one^\top w = 1
\]
has the unique solution $w_i = \softt_t(z_i + \nu^\star)$, $i = 1,\dots,N$, where $\nu^\star$
is a root of $h(\nu) := \sum_{i=1}^{N} \softt_t(z_i + \nu) - 1$.
\end{proposition}

The proof is in \cref{app:prox}. It runs in four moves: dualize the single affine constraint;
observe that the Lagrangian then separates across coordinates; solve each one-dimensional
subproblem in closed form, which is where soft thresholding appears; and recover $\nu^\star$ by
a one-dimensional root find, performed by bracketing and \textbf{bisection} to a tolerance of
$10^{-12}$. No closed form for $\nu^\star$ exists, since $h$ is piecewise linear with
data-dependent breakpoints.

Two properties then hold \emph{simultaneously}, which is the whole point of solving the prox
jointly: every coordinate with $|z_i + \nu^\star| \le t$ is set to \textbf{exactly zero} --- a
true floating-point zero, not a small number below a display threshold --- while
$\one^\top w = 1$ holds to the bisection tolerance. Neither is traded against the other.

\subsection{What this replaced: the prox-then-project bug}\label{sec:proxbug}

The first version of the solver did not compute the prox of $g$ jointly. It handled the two
nonsmooth pieces in sequence: soft-threshold first, then restore feasibility by a separate
Euclidean projection onto $\{\one^\top w = 1\}$, adding the uniform offset
$(1 - \one^\top w_{\text{half}})/N$ to \emph{every} coordinate. The failure is immediate once
stated: that offset is added to the coordinates soft thresholding had just set to zero,
reviving them. The returned portfolio satisfies the budget constraint and looks plausible, but
it is not sparse and it is not the minimizer.

The effect was not marginal. Against the independent solver of \cref{sec:insample}, on real
data with identical parameters, the two-step version reported \textbf{48 active positions where
the reference solver reported 9}, with an objective gap of \textbf{2.7--6\%}.

The general statement is that $\prox_{f_1 + f_2} \neq \prox_{f_1} \circ \prox_{f_2}$: the
proximal operator of a \emph{sum} of nonsmooth functions is not the composition of the
individual operators, and the two do not commute. When an objective carries several nonsmooth
components at once --- here an $\ell_1$ penalty and an affine constraint --- the prox of the sum
must be solved as one problem, which is what \cref{prop:jointprox} does.

Two points are worth drawing out. Each individual step in the broken version was correct in
isolation: soft thresholding is the right prox for $\ell_1$, and the uniform offset is the right
projection onto the hyperplane. Only their composition was wrong, which is why the error
survived inspection. And the output gave no signal --- a portfolio of 48 positions summing to
one is not visibly defective. The bug was found only by comparison against an independently
implemented solver, which is the argument for \cref{sec:insample} existing at all.

\subsection{Best-iterate tracking and the long-only variant}

A subgradient method is not a descent method: $f(w_{k+1}) \le f(w_k)$ is not guaranteed, and
the objective oscillates, particularly while $\alpha_k$ is large. The implementation therefore
records $\min_{j \le k} f(w_j)$ with the iterate attaining it and returns that, not the last
iterate, which would report an objective worse than the method achieved. Termination is judged
on the same quantity: the loop stops once the relative change in $\min_{j\le k} f(w_j)$ stays
below $10^{-8}$ for $100$ consecutive iterations.

For the backtest of \cref{sec:backtest}, $w \ge 0$ is added and the feasible set becomes the
probability simplex. Two things change. The proximal step becomes Euclidean projection onto
that simplex, computed by the standard sort-and-threshold algorithm \citep{duchi2008}. And the
$\ell_1$ term disappears: on the simplex $\norm{w}_1 = \one^\top w = 1$ identically, so
$\lambda\norm{w}_1$ contributes the constant $\lambda$ regardless of which feasible $w$ is
chosen and cannot influence the argmin. Hence $\lambda$ is dropped from the hyperparameter grid
in \cref{sec:backtest}: searching over it would consume compute and change nothing. The
subgradient \cref{eq:subgrad} is reused unchanged. The complete algorithm is stated as
pseudocode in \cref{alg:proxsub}.
